\chapter{КОМАНДЫ ВОСПРОИЗВЕДЕНИЯ}
\label{app:reproducibility}

Проверки выполняются из корня репозитория. Для generation core:

\begin{lstlisting}[language=bash]
cd generation_core
uv sync --frozen
uv run pytest
uv run ruff check .
uv lock --check
\end{lstlisting}

Для backend и frontend:

\begin{lstlisting}[language=bash]
cd backend
uv run pytest
uv run ruff check .
uv lock --check

cd ../frontend
bun run test
bun run lint
bun run build
bun run test:e2e
\end{lstlisting}

Внешний corpus не входит в репозиторий как vendored dataset. Его источник и commit фиксируются:

\begin{lstlisting}[language=bash]
git clone https://github.com/SDXorg/test-models.git out/sdxorg-test-models
git -C out/sdxorg-test-models checkout 21aab02739dc5187bc9564e4d3de14e575905d2f

uv run python external_benchmark.py \
  --root out/sdxorg-test-models \
  --output out/sdx-vensim.json
\end{lstlisting}

Пример benchmark-запуска локальной 8B после настройки OpenAI-compatible endpoint:

\begin{lstlisting}[language=bash]
export LOCAL_LLM_BASE_URL=http://127.0.0.1:18000/v1
export LOCAL_LLM_MODEL=Qwen/Qwen3-8B-GGUF
export LOCAL_LLM_STRUCTURED_OUTPUT=1
export LOCAL_LLM_MAX_TOKENS=8192

uv run python benchmark.py \
  --provider local \
  --prompts-file pipeline/benchmark_acceptance2_prompts.json \
  --runs 3 \
  --concurrency 1 \
  --max-repairs 2 \
  --output out/gpu-qwen3-8b-q4-v2-acceptance2.json \
  --artifacts-dir out/gpu-qwen3-8b-q4-v2-acceptance2
\end{lstlisting}

Для удалённой 8B используется Ministral-8B-2512 через официальный Mistral endpoint OpenRouter. Quantization endpoint не раскрывает, поэтому оно не предполагается. GPU и 14B в этой серии не используются. Секретный ключ читается из локального `.env` и не включается в отчёты:

\begin{lstlisting}[language=bash]
export OPENROUTER_MODEL=mistralai/ministral-8b-2512
export OPENROUTER_PROVIDER_ORDER=Mistral
export OPENROUTER_ALLOW_PROVIDER_FALLBACK=0
export LLM_SEED=2026
export LLM_STRUCTURED_OUTPUT=1
export LLM_COMPONENT_STRUCTURED_OUTPUT=1
export LLM_PATTERN_SHORTCUT=1
export LLM_ALLOW_UNSTRUCTURED_FALLBACK=0
export LLM_RAG_ENABLED=0

uv run --env-file ../.env python benchmark.py \
  --provider openrouter --mode ablation \
  --prompts-file pipeline/benchmark_confirmatory_8b_v2_prompts.json \
  --runs 3 --concurrency 1 --max-repairs 2 --corpus-role blind \
  --bootstrap-samples 10000 --seed 2026 \
  --output out/openrouter-ministral8b-ablation-v2.json \
  --artifacts-dir out/openrouter-ministral8b-ablation-v2-artifacts \
  --annotations out/openrouter-ministral8b-annotations-v2
\end{lstlisting}

Проверка controlled ablation на воспроизводимом fixture без обращения к LLM API:

\begin{lstlisting}[language=bash]
uv run python benchmark.py --provider fixture \
  --fixture tests/fixtures/llm_ablation_responses.json --mode ablation \
  --prompt-ids predator_prey_ru --corpus-role regression \
  --bootstrap-samples 200 --seed 19 \
  --artifacts-dir out/cli-ablation-artifacts \
  --annotations out/cli-annotations --output out/cli-ablation.json

uv run python evaluate.py compare --report out/cli-ablation.json \
  --bootstrap-samples 200 --seed 19 --output out/cli-comparison.json

uv run python evaluate.py agreement --first expert-1.csv \
  --second expert-2.csv --output out/annotation-agreement.json
\end{lstlisting}

Fixture проверяет только evaluation framework, а не качество LLM. Confirmatory run требует новый blind corpus и зафиксированный \codeid{pipeline/evaluation_protocol.json}. Эксперту передаются \codeid{expert/}, но не \codeid{private/answer_key.json}.

Отчёты содержат hashes, настройки и case results. Для публикации сохраняются версия vLLM plugin, inference lockfiles, prompts и аннотационная инструкция. Секреты в отчёты и репозиторий не включаются.

\section{Controlled локальные 8B и 30B-A3B}

Для совпадения pipeline hash используется исходник:
\begin{center}
\scriptsize\ttfamily d7ee55d505d0f3bb360ae5adc4156e956ff6cee9
\end{center}
Последующие изменения offline-анализа не подменялись на удалённой машине во время прогонов. Сначала отдельно проверена 8B, затем при неизменном source выполнена follow-up серия только 30B-A3B. Checkpoints и revisions:
\begin{center}
\scriptsize\ttfamily Qwen/Qwen3-8B-GGUF: 7c41481f57cb95916b40956ab2f0b139b296d974
\end{center}
\begin{center}
\scriptsize\ttfamily Qwen/Qwen3-30B-A3B-GGUF: e4d4bafdfb96a411a163846265362aceb0b9c63a
\end{center}

Запуск llama.cpp для 8B:
\begin{lstlisting}[language=bash,breakatwhitespace=false]
docker run --rm --gpus all -p 127.0.0.1:8000:8000 \
  -v "$HOME/.cache/huggingface:/root/.cache/huggingface" \
  ghcr.io/ggml-org/llama.cpp:server-cuda@sha256:5bb779a449a18d8f4333e9814e887746342fe555c528f5eabc92b2093e32673a \
  --model /root/.cache/huggingface/hub/models--Qwen--Qwen3-8B-GGUF/snapshots/7c41481f57cb95916b40956ab2f0b139b296d974/Qwen3-8B-Q4_K_M.gguf \
  --alias small --ctx-size 8192 --n-gpu-layers 99 \
  --host 0.0.0.0 --port 8000
\end{lstlisting}

Для 30B-A3B команда та же, за исключением пути и alias:
\begin{lstlisting}[language=bash,breakatwhitespace=false]
  --model /root/.cache/huggingface/hub/models--Qwen--Qwen3-30B-A3B-GGUF/snapshots/e4d4bafdfb96a411a163846265362aceb0b9c63a/Qwen3-30B-A3B-Q4_K_M.gguf \
  --alias large --ctx-size 8192 --n-gpu-layers 99
\end{lstlisting}

После готовности \codeid{/health} и сохранения \codeid{/props} выполняется benchmark без RAG и unstructured fallback. Для 8B задаётся \codeid{LOCAL_LLM_MODEL=small}, для 30B-A3B --- \codeid{large}; остальные настройки общие:
\begin{lstlisting}[language=bash]
export LOCAL_LLM_BASE_URL=http://127.0.0.1:8000/v1
export LOCAL_LLM_MODEL=large
export LOCAL_LLM_STRUCTURED_OUTPUT=1
export LOCAL_LLM_ALLOW_UNSTRUCTURED_FALLBACK=0
export LOCAL_LLM_MAX_TOKENS=4096
export LOCAL_LLM_TIMEOUT=600
export LOCAL_LLM_SEED=2026
export LLM_REQUEST_RETRIES=2
export LLM_COMPONENT_STRUCTURED_OUTPUT=0
export LLM_PATTERN_SHORTCUT=1
export LLM_RAG_ENABLED=0

uv run python benchmark.py --provider local --mode ablation \
  --prompts-file pipeline/benchmark_confirmatory_8b_v2_prompts.json \
  --runs 3 --concurrency 1 --max-repairs 2 --corpus-role blind \
  --bootstrap-samples 10000 --seed 2026 \
  --output out/local-qwen3-30b-a3b-q4km-ablation.json \
  --artifacts-dir out/local-qwen3-30b-a3b-q4km-artifacts \
  --annotations out/local-qwen3-30b-a3b-q4km-annotations

uv run python evaluate.py compare \
  --report out/local-qwen3-30b-a3b-q4km-ablation.json \
  --bootstrap-samples 10000 --seed 2026 \
  --output out/local-qwen3-30b-a3b-q4km-paired-analysis.json
\end{lstlisting}

Для 8B аналогично использованы \codeid{small} и префикс \codeid{local-qwen3-8b-q4km}. Парное сравнение моделей выполняется локально без GPU:
\begin{lstlisting}[language=bash]
uv run python evaluate.py compare-models \
  --baseline out/local-qwen3-8b-q4km-ablation.json \
  --treatment out/local-qwen3-30b-a3b-q4km-ablation.json \
  --bootstrap-samples 10000 --seed 2026 \
  --output out/local-qwen3-model-comparison.json
\end{lstlisting}

Фактические манифесты и полная статистика:
\begin{itemize}
    \item \codeid{evaluation/controlled-local-qwen3.json};
    \item \codeid{evaluation/controlled-local-qwen3-30b.json};
    \item \codeid{evaluation/controlled-local-qwen3-}\newline
      \codeid{paired-analysis.json};
    \item \codeid{evaluation/controlled-local-qwen3-}\newline
      \codeid{30b-paired-analysis.json};
    \item \codeid{evaluation/controlled-local-}\newline
      \codeid{qwen3-model-comparison.json}.
\end{itemize}

Raw artifacts и пакеты аннотации сохраняются локально в \codeid{generation_core/out/} и не коммитятся. Полный архив владельца --- \codeid{out/archives/sdcloud-controlled-qwen3-0.6.2.tar.gz}; его SHA-256 находится в \codeid{evaluation/controlled-local-qwen3-artifacts.sha256}. Он содержит приватные answer keys и не передаётся экспертам.

Для двух независимых экспертов подготовлен отдельный \codeid{out/archives/sdcloud-expert-review-0.6.2.tar.gz}: 432 задания и 251 model artifact в анонимизированных группах A/B, без answer keys и соответствия групп checkpoints. Приватное соответствие сохранено только в \codeid{out/archives/sdcloud-expert-review-private-0.6.2.tar.gz}. Контрольные суммы обоих архивов находятся в \codeid{evaluation/expert-review-packages.sha256}. После получения двух CSV выполняется \codeid{evaluate.py agreement} отдельно по группам A/B, а затем adjudication.

\codeid{compare-models} отклоняет несовпадение corpus, pipeline, protocol, LLM settings и prompt/run/seed пар. После каждой серии evaluation inference-сервер и GPU-monitor останавливаются; после 30B-A3B временные контейнеры удалены. Основной продуктовый стек при этом не используется и не останавливается.

\section{Полный REST product E2E}

Проверка использует отдельный local/test-стек. В нём доступен одноразовый код для разработки, а генерация заменена воспроизводимым fixture. Production и внешний LLM API не вызываются:

\begin{lstlisting}[language=bash]
cd backend
uv run python tests/live_e2e.py \
  --base-url http://127.0.0.1:61231 \
  --output .test_tmp/rest-e2e-report.json
\end{lstlisting}

Сценарий проверяет OTP и refresh, создание беседы, три упорядоченных сообщения и идемпотентный replay, контекстный ответ без модели, вторую версию после правки, симуляцию, Python/Vensim/XMILE exports и запрет доступа другого пользователя. Фактический изолированный прогон завершился 46 HTTP-шагами и 28 успешными проверками; его volumes и контейнеры удалены после теста.
